\scanpage{47}
% archive sheet number: 46

\textcircled{17} \uline{Anneau parfait de car. $p$} ($p>0$

\[
 p\,1_A=0,\qquad A\xrightarrow{x\mapsto x^p}A\quad\text{isom}
\]

% slanted
\margin{en car. $0$, on convient que parfait signifie réduit ?}

NB\quad parfait $\Longrightarrow$ réduit

\[
 X_{p\text{-perf}}=\widetilde{R_{\unsure{\mathrm{top}|\mathrm{sch}}\,p\text{ perf}}}
 =\widehat{R_{p\text{-perf}}}
\]

$R_{p\text{ perf}}$ catégorie des schémas affines \deleted{\unsure{plats|de}} car. $p$ parfaits qui sont clôtures parfaites de schémas de type fini sur $\mathbb Z/p\mathbb Z$.

top. $p$-perf sur $R$ : $X_i\longrightarrow X$ est couvrante si $\exists i_0$, $(X_{i_0})_p\deleted{\ensuremath{\unsure{\times|\to}}}\to(X)_p$ admet une section après passage à la clôture parfaite de $X_p$

\textcircled{18} \uline{\deleted{\unsure{Anneau de corps|Anneau de car.}} \inserted{compact} parfait de car. $p$} ($>0$).

parfait de car. $p$ $\cap$ pr. corps

\[
\begin{aligned}
 X_{\deleted{\ensuremath{\unsure{\mathrm{pr\,corps}|\mathrm{absol.}}}}\,p\text{-perf}}
 &\simeq \widetilde{R_{\deleted{\ensuremath{\unsure{\mathrm{perf}|\mathrm{pr}}}}\,\mathrm{pr\,corps}\,p\text{-perf}}}.\\
 &\simeq R_{\mathrm{cons}\,p\text{-perf}}
\end{aligned}
\]

top. \unsure{quelq. p. parfaite|quasi p parfaite|cpct p-parfaite} : $X_i\to X$ est couv\textsuperscript{te} ssi $\exists i_0$ t.q. $X_{i_0\,p}\longrightarrow X_p$ a un pt radiciel au dessus de tt pt de $X_p$

$R_{\mathrm{cons}\,p\text{-perf}}=$ catégorie des clôtures parfaites des schémas de car. $p$ dans $R_{\mathrm{cons}}$.

\textcircled{19} \uline{Corps parfait} \deleted{\unsure{$(p>0)$|de $p$}} \inserted{(car $p>0$)}

\[
\begin{aligned}
 X_{\mathrm{corps}\,p\text{-perf}}&\simeq
 \widetilde{R_{\mathrm{top\,corps}\,p\text{-perf}}}\simeq\deleted{\ensuremath{\unsure{R|P}}}\\[-2pt]
 &\hspace{1em}\widetilde{R_{\unsure{\mathrm{cons}_p|\mathrm{cons}/p}\,\deleted{\ensuremath{\mathrm{perf}}}\,\mathrm{Zar}}}
\end{aligned}
\]

$X_i\longrightarrow X$ est couvr\textsuperscript{te} ssi $X_{i\,p}\to X_p$ a des pt radiciels \unsure{pts un|fts un|pts au|sur chaque} par fibre.

\deleted{\textcircled{20} \uline{Anneau à caractéristique}}

\deleted{(Anneau parfait (ou compact parfait, ou corps parfait) de car. $p>0$ non précisé (inclut car $0$)}

\deleted{\unsure{$\Delta$.|$\triangle$.|$A$.} On veut que loc sur $\mathbb Z$ il ait les propriétés envisagées.}

\deleted{Soit $C$ l'ens des nbs qui sont soit premiers, soit $0$ ($0$ mis en cas de car. $0$), \deleted{\unsure{et on|donc|Donc}} soit le symbole $\varnothing$. Alors $\mathcal E$. se décompose can\textsuperscript{ment} en $\mathcal E\simeq\coprod_{c\in C}\mathcal U_c$, avec $\mathcal E$ $\unsure{\Delta|A}\mapsto\left.\vphantom{U}\right|\mathcal U_c$ de car. $p$ si $c\ne\varnothing$, et $c\in C$, \unsure{$\Delta|\mathcal U_\varnothing$|$\Delta\mathcal U_\varnothing$}$=0$. Soient alors $X_{p\deleted{\ensuremath{\mathrm{perf}}}}\subset X_{\mathrm{ann}}$, s-to\unsure{pos|pes} immersions}

\deleted{$X_{\infty\deleted{\ensuremath{\unsure{\mathrm{perf}|\mathrm{cons}}}}}\subset$}
